\documentclass{ieeeaccess}
\usepackage{cite}
\usepackage{amsmath,amssymb,amsfonts}
\usepackage{algorithm}
\usepackage{algorithmic}
\usepackage{graphicx}
\usepackage{textcomp}
\usepackage{booktabs}
\usepackage{multirow}
\usepackage{url}
\def\BibTeX{{\rm B\kern-.05em{\sc i\kern-.025em b}\kern-.08em
    T\kern-.1667em\lower.7ex\hbox{E}\kern-.125emX}}
\begin{document}
\history{Submitted on May 14, 2026.}
\doi{10.1109/ACCESS.XXXX.DOI}

\title{Precision-Aware Residency Scheduling for Mixed-Precision Large Language Model Decoders}
\author{\uppercase{Author One}\authorrefmark{1}, \uppercase{Author Two}\authorrefmark{1},
        and \uppercase{Author Three}\authorrefmark{1}}
\address[1]{Department of Computer Science and Engineering, Ewha Womans
University, Seoul, Republic of Korea (corresponding author e-mail:
{authorone}@ewha.ac.kr)}
\tfootnote{This work was supported in part by [funding agency and grant number TBD].
The code, configurations, and the results that back every figure and table in
this article are available in the project repository accompanying the submission.}

\markboth
{Author~\headeretal: Precision-Aware Residency Scheduling for Mixed-Precision LLM Decoders}
{Author~\headeretal: Precision-Aware Residency Scheduling for Mixed-Precision LLM Decoders}

\corresp{Corresponding author: Author One (e-mail: authorone@ewha.ac.kr).}

\begin{abstract}
Autoregressive decoding in large language models is bound by per-token
weight traffic from off-chip memory: every weight is read once per
generated token, so latency and energy track the bytes transferred
rather than floating-point operations. Layer-wise mixed precision
shrinks those bytes, but a deployment-time scheduler that does not
know each layer's bit width must budget every weight tile at the
worst-case width, leaving on-chip residency capacity unused. We
formalize the resulting decision as precision-aware residency
scheduling: given a fixed mixed-precision policy and an on-chip
buffer, decide which weight tiles to pin so that per-token traffic to
dynamic random-access memory is minimized. We instantiate the
scheduler as a greedy zero-one knapsack on per-tile bytes and compare
it against a precision-oblivious baseline that uses eight-bit budgets.
Across four open-weight checkpoints spanning approximately 1.5, 2.6,
7.6, and 8.0 billion parameters and three model families (Qwen, Gemma,
Llama), the precision-aware scheduler reduces per-token off-chip
traffic by 1.33 times to 1.97 times at a comparable operating point
near half the four-bit model footprint, and up to 3.17 times at the
knee where the buffer matches the four-bit footprint of the smallest
model. The mixed-precision policy commonly used as input
survives joint post-training quantization with a Kendall rank
correlation of $+$1.0 against the predicted ordering. A
graphics-processing-unit profile of the unquantized baseline
corroborates the analytic ranking. All code, configurations, and
results are released for reproducibility.
\end{abstract}

\begin{keywords}
Large language model inference, mixed-precision quantization,
on-chip residency, off-chip memory traffic, knapsack scheduling,
hardware--software co-design, post-training quantization, hardware
accelerators.
\end{keywords}

\titlepgskip=-15pt

\maketitle

\section{Introduction}
\label{sec:introduction}
\PARstart{L}{arge} language models (LLMs) generate text autoregressively:
each new token requires a forward pass through the full decoder stack
with the keys and values of all prior tokens cached on chip and a fresh
read of the weight matrices from off-chip memory. In the decode regime,
the batch size is one, the arithmetic intensity is low, and the
processing-element (PE) array of the host graphics processing unit (GPU)
or accelerator spends most of its cycles waiting for memory traffic
\cite{sheng2023flexgen}. Reducing the per-token weight bytes is
therefore the dominant lever on both latency and energy.

Layer-wise mixed-precision quantization is one of the canonical
techniques for this reduction. Hessian-aware quantization (HAWQ)
\cite{dong2019hawq} and its eigenvalue variant HAWQ-V2
\cite{dong2020hawqv2} assign different bit widths per layer guided by
the loss landscape; block-wise reconstruction error (BRECQ)
\cite{li2021brecq} formulates the same target with a different signal;
generalized post-training quantization (GPTQ) \cite{frantar2023gptq}
makes the quantization itself Hessian-aware. Recent LLM accelerator
designs --- including our own laboratory's prior power-aware mapping
work and LLM-specific accelerator (both anonymized for review) ---
take mixed-precision policies as system input and search the dataflow
and mapping space to extract energy efficiency on top.

This paper asks a question that sits between the algorithm and the
mapping layers: \emph{what should the deployment-time scheduler do with
the mixed-precision policy at runtime?} Concretely, when the on-chip
buffer is too small to hold the full decoder, which weight tiles do we
keep resident so that subsequent tokens reuse them without re-streaming
from dynamic random-access memory (DRAM)?

We make three observations.
\begin{itemize}
\item \textbf{Observation~1.} In the autoregressive decode regime,
search of the loop-mapping space does not break the byte-traffic
monotone. We show empirically in Section~\ref{sec:motivation} that the
optimal mapping for the same general matrix multiplication (GEMM) at
four-bit and at eight-bit differ in tile geometry, but the cycle count
and energy in a representative 16$\times$16 PE array track byte width
almost exactly, so the mapping contribution is absorbed by the byte
cost.
\item \textbf{Observation~2.} The natural surface for adding scheduling
intelligence is \emph{residency}. With the model fixed, every weight is
read exactly once per token. If a tile fits in the on-chip buffer, it
can be pinned once and reused across many tokens. The decision becomes
a classical zero-one knapsack with a single capacity constraint.
\item \textbf{Observation~3.} The mixed-precision policy changes the
knapsack instance, but naive deployment ignores it. A
precision-oblivious scheduler that budgets every tile at INT8 (the
worst-case bit width in our setting) pins fewer tiles than a
precision-aware scheduler that uses the actual per-tile bytes from the
policy. At fixed buffer size, the aware scheduler streams 1.3 to 3.2
times fewer bytes for the same accuracy.
\end{itemize}

\noindent\textbf{Contributions.} The contributions of this paper are
the following.
\begin{itemize}
\item We formulate precision-aware residency scheduling as a knapsack
on per-tile weight bytes and quantify its win against a
precision-oblivious INT8-budget baseline across three LLM families
(Qwen, Gemma, Llama) spanning four checkpoints from 1.5 to 8 billion
parameters.
\item The pattern is universal across the four checkpoints we sweep:
Qwen2.5-1.5B-Instruct, Gemma-2-2B-IT, Qwen2.5-7B-Instruct, and
Llama-3.1-8B-Instruct all exhibit a monotone ranking under joint
quantization and an aware-versus-oblivious win in the same direction;
the win factor grows with the buffer-to-decoder ratio, reaching 1.80
to 1.97 times for the seven- and eight-billion-parameter models at a
buffer of roughly 60\,\% of the four-bit decoder footprint.
\item As a method extension, replacing the per-layer Hessian-aware
score with a per-tile bits-per-byte score produces qualitatively
different policies that allocate the eight-bit budget to attention
key-value (KV) projections at the same average bit-width budget.
Real GPTQ verification across all four models (Qwen-1.5B/7B,
Gemma-2-2B, Llama-8B) reduces WikiText-2 perplexity in seven of eight
matched-budget settings, with the largest gain ($-0.162$ at
$\bar B\!=\!4.5$ on Qwen-1.5B) concentrated at the tightest budget
where protecting KV projections matters most.
\item We also separate the contribution of \emph{precision awareness}
from the choice of solver: replacing the smallest-first knapsack with
a uniformly-random ordering of the same byte-aware budget shifts the
DRAM traffic by less than one percent across all sixteen
(model, policy) settings. The win is therefore attributable to the
byte-awareness of the budget, not to a particular packing heuristic,
which makes the technique robust to deployment-specific solver
choices.
\item We provide a complete reproducibility trail, including a
sudo-less source build of Timeloop and Accelergy inside a single
conda environment, a documented gcc-13 patch needed to compile
ninety-two affected headers, and a deterministic driver that
regenerates every figure and table from the same set of pipeline
modules.
\end{itemize}

\noindent\textbf{Stack position.} The literature treats LLM efficiency
at three layers (Fig.~\ref{fig:stack}). Our contribution is the
scheduling layer; the upper two layers are inputs.

\begin{figure}[t]
\centering
\begin{tabular}{|l|}
\hline
\textsf{Algorithm layer (input).} \\
HAWQ \cite{dong2019hawq}, BRECQ \cite{li2021brecq}, GPTQ \cite{frantar2023gptq}. \\
Output: per-(layer, projection) bit widths. \\ \hline
\textsf{Mapping layer (input).} \\
Timeloop \cite{parashar2019timeloop}, MAGNETO and GATHER (anonymized). \\
Output: per-shape tile sizes, loop-nest order. \\ \hline
\textbf{\textsf{Scheduling layer (this paper).}} \\
Input: mixed-precision policy and on-chip buffer of size $B$. \\
Decision: which weight tiles to pin on-chip. \\
Output: per-token off-chip bytes (efficiency). \\ \hline
\end{tabular}
\caption{The three-layer stack. Mapping and scheduling are sometimes
conflated in single-pass dataflow papers; we show in
Section~\ref{sec:motivation} that under autoregressive decode the
mapping layer's degrees of freedom collapse onto byte traffic, so the
scheduling layer's contribution is additive to mixed precision rather
than redundant with it.}
\label{fig:stack}
\end{figure}

\section{Background}
\label{sec:background}

\subsection{Mixed-Precision Quantization and Policy Search}
Per-layer sensitivity scoring assigns each transformer block a score
that measures the perplexity (PPL) degradation caused by quantizing
that block to a specified number of bits. Common formulations use
Hessian-trace signals \cite{dong2019hawq, dong2020hawqv2}, block
reconstruction error \cite{li2021brecq}, or simple ablation. The
\emph{policy search} problem is then: given allowed widths
$W = \{b_1,\ldots,b_k\}$ and an average bits-per-weight (BPW) budget
$\bar{B}$, find a per-layer assignment $B[l] \in W$ that minimizes the
total $\Delta \mathrm{PPL}$ subject to $\mathrm{mean}(B) \le \bar{B}$.
Greedy promotion by $\Delta \mathrm{PPL}$ per added bit is the standard
solver and is near-optimal under monotone scores. Formally, the greedy
rule promotes the layer with the largest gain in $\Delta \mathrm{PPL}$
per added bit at each step:
\begin{equation}
\mathrm{score}(l, b\!\to\!b') = \frac{\Delta\mathrm{PPL}(l,b) - \Delta\mathrm{PPL}(l,b')}{b' - b}.
\label{eq:greedy_score}
\end{equation}

\subsection{Post-Training Quantization with GPTQ}
GPTQ \cite{frantar2023gptq} performs greedy column-by-column
quantization while updating the remaining weights to compensate for
accumulated error, using an approximate inverse Hessian computed from a
small calibration set. It absorbs a meaningful fraction of the loss
that isolated-layer sensitivity predicts, as we quantify in
Section~\ref{sec:e1}.

\subsection{Dataflow Mapping for Matrix Multiplication}
Timeloop \cite{parashar2019timeloop} searches the loop-nest tiling
space for a given problem ($M\!\times\!K\!\times\!N$) on a given
accelerator (PE array, memory hierarchy, bandwidth). For an
Eyeriss-like architecture \cite{chen2017eyeriss} with a 16$\times$16 PE
array, a 256-KB on-chip ``weight-input buffer,'' and a 24-bit
``accumulation buffer,'' Timeloop reports cycles, energy from its
built-in power--area--timing model, and DRAM byte traffic per GEMM.

\subsection{LLM Accelerator Memory Hierarchy}
Modern LLM accelerators are designed around the observation that
weights dominate memory traffic. Common patterns include large on-chip
static random-access memory (SRAM) banks, or high-bandwidth memory
(HBM) used as an explicit cache tier, that hold a fraction of the model
resident across token boundaries
\cite{lu2021sanger, guo2022ant}. This paper focuses on that residency
layer: given the policy and the buffer size, which tiles stay.

\section{Motivation: Mapping Search Does Not Save Decode}
\label{sec:motivation}

We instantiate the precision-aware-mapping hypothesis as follows. For
each unique linear shape $s$ in a Qwen2.5-1.5B decoder block (seven
shapes: \texttt{q\_proj}, \texttt{k\_proj}, \texttt{v\_proj},
\texttt{o\_proj} and \texttt{gate\_proj}, \texttt{up\_proj},
\texttt{down\_proj}) and each bit width $b \in \{4, 8\}$, we run the
Timeloop mapper on a 16$\times$16 PE array with a 256-KB on-chip
buffer and report cycles, energy, DRAM bytes, and PE utilization. We
then compute the cross-application loss
\begin{equation}
L_{a \to b} := \frac{\mathrm{metric}(\mathrm{best}_a,\ \mathrm{weights}\!=\!b)}{\mathrm{metric}(\mathrm{best}_b,\ \mathrm{weights}\!=\!b)},
\label{eq:cross_app}
\end{equation}
that is, how much we give up by forcing the wrong-bit-width's best
mapping onto the actual bit-width's data.

\paragraph{Result} The mappings differ: at INT4 the mapper picks a tile
that holds the entire output dimension in the global buffer, while at
INT8 the mapper must split the output dimension because the same tile
would overflow the buffer. However, the cycles and energy outcomes are
dominated by byte width
(Table~\ref{tab:mapping_pilot}).

\begin{table}[t]
\caption{Mapping Pilot on Qwen2.5-1.5B: Cycle/Energy Ratios Track Byte Width}
\label{tab:mapping_pilot}
\centering
\begin{tabular}{lrrrr}
\toprule
Shape & $b$ & Cycles & Energy (pJ) & DRAM (MB) \\
\midrule
\texttt{q\_proj}  & 4 & 147{,}552 & $1.35\!\times\!10^8$ & 1.18 \\
\texttt{q\_proj}  & 8 & 147{,}648 & $2.70\!\times\!10^8$ & 2.36 \\
\texttt{up\_proj} & 4 & 860{,}256 & $7.85\!\times\!10^8$ & 6.89 \\
\texttt{up\_proj} & 8 & 860{,}256 & $1.57\!\times\!10^9$ & 13.77 \\
\bottomrule
\end{tabular}
\end{table}

The cycle ratio is $\approx 1.000$; the energy and DRAM ratios are
$\approx 2.000$, the byte ratio of $b=8$ to $b=4$. The mapping geometry
difference does not translate into a meaningful cycle delta because at
$N\!=\!1$, every weight is touched once and the total DRAM time equals
total bytes divided by bandwidth, which is identical between the two
mappings: they stream the same total weight bytes.

\paragraph{Conclusion of Section~\ref{sec:motivation}} For
autoregressive decode on an Eyeriss-class array, mapping search is not
the right surface for adding scheduling intelligence on top of mixed
precision. The right surface is the layer above: which tiles do we keep
on-chip across tokens?

\section{Method: Precision-Aware Residency Scheduling}
\label{sec:method}

\subsection{Problem Definition}
Let a decoder consist of $L$ blocks, each containing seven projections
indexed as tiles $t \in T$. Each tile has a weight count $n_t$ (a
function of the model hidden size and the projection role). Given a
layer-wise mixed-precision policy $B[t] \in \{4, 8\}$, each tile has
bytes
\begin{equation}
\mathrm{bytes}_{\mathrm{aware}}(t) = \frac{n_t \cdot B[t]}{8}.
\label{eq:bytes_aware}
\end{equation}
Per token, the GEMM for each tile is read once. If the tile is
\emph{pinned} in the on-chip buffer of capacity $B_{\mathrm{GBuf}}$,
the read is free for that token and every subsequent token in the
sequence; if it is \emph{streamed}, it costs
$\mathrm{bytes}_{\mathrm{aware}}(t)$ of off-chip traffic per token.

The scheduler selects a subset $P \subseteq T$ such that
$\sum_{t \in P} \mathrm{budget}_{\mathrm{bytes}}(t) \le B_{\mathrm{GBuf}}$.
The three schedulers we compare differ in the byte function they use
for the constraint and in how they choose tiles within the budget:
\begin{itemize}
\item \textbf{Precision-aware (smallest-first knapsack):} uses
$\mathrm{bytes}_{\mathrm{aware}}(t)$ and visits tiles in ascending
order of size (Algorithm~\ref{alg:greedy}). Our default scheduler.
\item \textbf{Precision-aware (random order):} uses the same
$\mathrm{bytes}_{\mathrm{aware}}(t)$ for the constraint but visits
tiles in a uniformly-random order, averaged over eight seeds. Isolates
the contribution of the smallest-first rule from precision awareness.
\item \textbf{Precision-oblivious (INT8 budget):} uses
$\mathrm{bytes}_{\mathrm{int8}}(t) = n_t$ for the constraint while still
consuming the actual $\mathrm{bytes}_{\mathrm{aware}}(t)$ once a tile is
pinned. Reflects a deployment scheduler that ignores the policy and
pessimistically assumes the worst-case bit width.
\end{itemize}
All three schedulers maximize the bytes saved per token,
$\sum_{t \in P} \mathrm{bytes}_{\mathrm{aware}}(t)$, subject to their
own budget constraint. We will show in Section~\ref{sec:e3_sweep} that
the two precision-aware schedulers (smallest-first vs random) finish
within roughly one percent of each other --- the win over oblivious is
almost entirely attributable to byte-awareness, not to the
smallest-first rule, which makes the result robust to solver choice.

\subsection{Solver: Smallest-First Greedy Knapsack}
\label{sec:solver}
For autoregressive decode at $N\!=\!1$, every tile is accessed exactly
once per token. The access frequency does not vary by tile, so the
per-byte saving rate is identical for every tile. Under that
condition, the optimal zero-one knapsack reduces to greedy
smallest-first: sort tiles ascending by $\mathrm{budget}_{\mathrm{bytes}}$
and pin until the next one would overflow $B_{\mathrm{GBuf}}$. This is
optimal for both schedulers and runs in $O(|T|\log|T|)$ time
(Algorithm~\ref{alg:greedy}).

\begin{algorithm}[t]
\caption{Greedy Smallest-First Residency}
\label{alg:greedy}
\begin{algorithmic}[1]
\REQUIRE tiles $T$, budget function $f$, capacity $B_{\mathrm{GBuf}}$
\STATE sort $T$ ascending by $f(t)$
\STATE $P \gets \emptyset$;\ \ $u \gets 0$
\FORALL{$t \in T$ in sorted order}
  \IF{$u + f(t) \le B_{\mathrm{GBuf}}$}
    \STATE $P \gets P \cup \{t\}$;\ \ $u \gets u + f(t)$
  \ENDIF
\ENDFOR
\RETURN $P$
\end{algorithmic}
\end{algorithm}

\subsection{Hardware-Normalized Policy (Per-Tile, Optional)}
\label{sec:hwnorm}
The per-layer greedy of \eqref{eq:greedy_score} promotes layer $l$
based on $\Delta\mathrm{PPL}/\Delta\mathrm{bits}$, which is independent
of how many bytes the promotion costs. With a homogeneous decoder all
layers are the same size, so this is fine; but within a layer, the seven
projections have very different sizes
($\texttt{k\_proj}$ and $\texttt{v\_proj}$ are six times smaller than
$\texttt{gate\_proj}$ in Qwen2.5-1.5B). A per-tile hardware-normalized
score uses $\Delta\mathrm{PPL}$ per added byte,
\begin{equation}
\mathrm{score}_{\mathrm{hwnorm}}(t, b\!\to\!b') = \frac{\Delta\mathrm{PPL}(t,b) - \Delta\mathrm{PPL}(t,b')}{n_t \cdot (b'-b)/8},
\label{eq:hwnorm}
\end{equation}
which promotes the smallest highly sensitive tiles first (typically
$\texttt{k\_proj}$ and $\texttt{v\_proj}$ of attention) and leaves the
large multi-layer-perceptron (MLP) projections at the minimum width. We
measure per-tile $\Delta\mathrm{PPL}$ by ablating each of the
$L\!\times\!7$ linear layers independently.

\section{Experimental Setup}
\label{sec:setup}

\subsection{Models, Weights, and Calibration}
We evaluate on four open-weight chat-tuned LLM checkpoints spanning
three families (Qwen, Gemma, Llama) and four scales from 1.54 to 8.03
billion parameters, listed in Table~\ref{tab:models}. All baselines are
fetched and stored read-only; the experimental pipeline produces
quantized copies in a separate directory and never modifies the
baseline weights. For Llama-3.1-8B-Instruct we use the
\texttt{NousResearch/Meta-Llama-3.1-8B-Instruct} community
redistribution, which carries the same Git SHA as Meta's gated repo;
the two are bit-identical and we treat them as the same checkpoint for
the purposes of this paper.

\begin{table}[t]
\caption{Models and Decoder Footprints}
\label{tab:models}
\centering
\begin{tabular}{lrrr}
\toprule
Model & Params & INT4 dec.\ & INT8 dec.\ \\
\midrule
Qwen2.5-1.5B-Instruct      & 1.54\,B & 655\,MB & 1.31\,GB \\
Gemma-2-2B-IT              & 2.61\,B & 1.04\,GB & 2.07\,GB \\
Qwen2.5-7B-Instruct        & 7.61\,B & 3.25\,GB & 6.50\,GB \\
Llama-3.1-8B-Instruct$^*$  & 8.03\,B & 3.49\,GB & 6.98\,GB \\
\bottomrule
\multicolumn{4}{l}{\footnotesize $^*$ obtained from the
\texttt{NousResearch/Meta-Llama-3.1-8B-Instruct}}\\
\multicolumn{4}{l}{\footnotesize \phantom{$^*$} community redistribution
at the same Git SHA as Meta's gated repo.}
\end{tabular}
\end{table}

Sensitivity scoring ablates one decoder block, or one projection in the
per-tile variant of Section~\ref{sec:hwnorm}, to bits
$b \in \{4, 8\}$ using symmetric per-output-channel uniform
fake-quantization, and reports WikiText-2 test-set perplexity over a
sliding-window of 2048 tokens with 200 calibration texts. For
Qwen2.5-7B and Llama-3.1-8B the models are sharded across the
available GPUs via Hugging Face Accelerate \cite{accelerate} (using
\texttt{device\_map=\"auto\"}) with 50 calibration texts for the 7B and
the default 200 for the 8B (which fits in four 16-GiB cards in
half-precision); the resulting $\Delta\mathrm{PPL}$ rankings are
consistent with the 200-text setting used for the 1.5B model.

GPTQ uses 32 calibration samples from WikiText-2 train with a group
size of 128 and the descending-activation order option enabled. For
mixed-precision policies, we use a regular-expression override on the
quantizer configuration to set eight-bit on the layers selected by the
policy. For the 7B model we pass \texttt{device\_map=\"auto\"} to enable
multi-GPU-sharded calibration.

Evaluation reports WikiText-2 test perplexity over a sliding window of
2048 tokens on the full split. GPTQ checkpoints are loaded through the
quantizer's own loader rather than the generic
\texttt{AutoModelForCausalLM} loader to avoid an Optimum-Transformers
compatibility issue in the installed library versions.

\subsection{Hardware Model}
Timeloop v3.0.3 with the necessary patches drives the mapping pilot of
Section~\ref{sec:motivation}. The architecture template is a
16$\times$16 PE array with a 256-KB on-chip ``weight-input buffer'' and
a DRAM bandwidth of 32~GB/s. Per-bit multiply-accumulate (MAC) and
storage energy come from Timeloop's built-in power--area--timing model
with the word-width parameter set to the operand bit width.

The residency experiments of Sections~\ref{sec:e2_residency}--%
\ref{sec:e4} use an analytic byte-traffic model: every weight is
touched exactly once per token, so DRAM traffic per token equals
$\sum_{t \notin P} \mathrm{bytes}_{\mathrm{aware}}(t)$. The negative
mapping result of Section~\ref{sec:motivation} justifies this: at
$N\!=\!1$, total cycles track DRAM bytes linearly and the mapping
geometry contribution is below noise.

\paragraph{Energy per byte}
We translate per-token DRAM bytes into a Joule estimate by harvesting
the energy-per-byte ratio from the same Timeloop runs that drive the
mapping pilot. Averaging across the six (shape $\times$ bit-width)
mapper runs we observe $\mathrm{energy/byte}\!=\!114.03$\,pJ at INT4 and
$114.14$\,pJ at INT8 --- statistically identical (ratio $1.001$), in
direct correspondence with the negative mapping result of
Section~\ref{sec:motivation}. We therefore take a single constant
$E_b\!=\!114.08$\,pJ/byte and report energy per token as $E_b$ times
the DRAM-byte traffic reported by the residency model. This gives a
defensible energy lower bound that ignores PE-array idle power and
on-chip-buffer refresh, both of which are common to all schedulers.

\subsection{GPU Profile (Sanity Check)}
We separately measure GPU-side latency, throughput, and approximate
energy per token on a single NVIDIA RTX 2000 Ada (16 GiB) using the
NVIDIA Management Library (NVML) for power sampling at 20-ms intervals
through a background thread. These numbers serve as a sanity check for
the analytic model and as a deployment-scale reference; they are not
the main result.

\section{Results}
\label{sec:results}

\subsection{E1: Ranking Validation}
\label{sec:e1}
We measure four policies per model: INT4-uniform, INT8-uniform, and
two mixed policies from the greedy per-layer score at average bits per
weight $N \in \{4.5, 5.0\}$. For each we quantize with GPTQ and measure
real WikiText-2 perplexity.

\begin{table}[t]
\caption{Real GPTQ Perplexity Across Four Models.\\ All Four Are Monotone in Average Bits per Weight (Kendall $\tau = -1.0$)}
\label{tab:e1}
\centering
\begin{tabular}{lrrrr}
\toprule
Policy & Qwen-1.5B & Gemma-2B & Qwen-7B & Llama-8B \\
\midrule
FP16 baseline & 8.890  & ---     & 6.121$^*$ & ---     \\
INT4-uniform  & 10.452 & 14.800  & 7.785     & 7.873   \\
Mixed-4.5     & 10.279 & 14.723  & 7.724     & 7.561   \\
Mixed-5.0     & 10.091 & 14.387  & 7.661     & 7.514   \\
INT8-uniform  & 9.604  & 13.773  & 7.376     & 7.159   \\
\bottomrule
\multicolumn{5}{l}{\footnotesize $^*$ measured with 50 calibration
texts, multi-GPU sharded.}
\end{tabular}
\end{table}

All four models exhibit a perfectly monotone ranking (Table~\ref{tab:e1});
the Kendall rank correlation between average bits per weight and
perplexity is $-1.0$ for every model.

\paragraph{Sensitivity is conservative}
For Qwen-1.5B, the isolated-layer upper bound predicts a
$\Delta \mathrm{PPL}$ of $+$4.00 for INT4-uniform versus FP16; the real
GPTQ value is $+$1.56, that is, GPTQ absorbs 61\% of the predicted loss
through Hessian-aware reconstruction. The ratio is similar for Gemma-2B,
Qwen-7B, and Llama-8B. The \emph{ranking} of policies, however, matches the
upper-bound ranking exactly (Kendall $\tau$ on the real-versus-upper
sequence is $+$1.000 for Qwen-1.5B; see
Fig.~\ref{fig:e1_ranking_validation}). The implication is that
sensitivity is an order oracle, not a magnitude oracle.

\paragraph{GPTQ floor}
Even uniform-INT8 GPTQ adds $+$0.71 perplexity points on Qwen-1.5B
versus the FP16 baseline, against a sensitivity upper bound of only
$+$0.03. This is a property of GPTQ's group-wise calibration with
group size 128 and 32 calibration samples and would shrink under larger
calibration; for relative policy comparison this offset cancels out.

\paragraph{Downstream-task validation}
To check that the perplexity ranking carries over to a downstream task,
we evaluate zero-shot HellaSwag accuracy on a 1000-sample validation
subset (standard error $\approx$ 1.5 percentage points at this $N$)
using length-normalized log-likelihood scoring
(Table~\ref{tab:hellaswag}). All sixteen quantized configurations are
within three percentage points of their FP16 baseline. The largest
absolute degradation is on Llama-8B INT4-uniform ($-2.8$\,pp versus
FP16), and the Mixed-4.5 policy recovers it to within $-1.4$\,pp at a
matched average bit-width below five bits. On the smaller models the
INT4 perplexity gap does not produce a statistically detectable
HellaSwag gap, consistent with the coarse-grained nature of four-way
multiple-choice accuracy versus token-level perplexity.

\begin{table}[t]
\caption{Zero-Shot HellaSwag Accuracy (1000-Sample Subset): Quantization Is Safe for Downstream Task Across All Four Models}
\label{tab:hellaswag}
\centering
\small
\begin{tabular}{lrrrrr}
\toprule
Model      & FP16  & INT4-unif & Mixed-4.5 & Mixed-5.0 & INT8-unif \\
\midrule
Qwen-1.5B  & 0.573 & 0.564 & 0.569 & 0.560 & 0.568 \\
Gemma-2-2B & 0.419 & 0.442 & 0.428 & 0.422 & 0.422 \\
Qwen-7B    & 0.672 & 0.671 & 0.666 & 0.669 & 0.673 \\
Llama-8B   & 0.677 & 0.649 & 0.663 & 0.663 & 0.673 \\
\bottomrule
\end{tabular}
\end{table}

\begin{figure}[t]
\centering
\includegraphics[width=\columnwidth]{../experiments/results/figs/e1_ranking_validation.png}
\caption{E1 ranking validation for Qwen2.5-1.5B-Instruct. Left:
sensitivity upper bound (dashed) versus real GPTQ perplexity (solid);
right: GPTQ absorbs most of the predicted loss, but the ranking is
preserved.}
\label{fig:e1_ranking_validation}
\end{figure}

\subsection{E2-Residency: Main Result}
\label{sec:e2_residency}

\begin{figure}[t]
\centering
\includegraphics[width=\columnwidth]{../experiments/results/figs/codesign_main.png}
\caption{Accuracy versus per-token off-chip traffic at
$B_{\mathrm{GBuf}}\!=\!512$\,MB on Qwen2.5-1.5B. Filled circles are
precision-aware schedulers; crosses are precision-oblivious schedulers.
The arrows show the per-policy reduction in off-chip traffic from using
the precision-aware scheduler.}
\label{fig:codesign_main}
\end{figure}

Figure~\ref{fig:codesign_main} is our main result for Qwen-1.5B at the
$B_{\mathrm{GBuf}}\!=\!512$\,MB operating point, slightly below the
INT4-uniform 655-MB decoder total. Filled circles are the
precision-aware schedulers; crosses are the precision-oblivious INT8
budget schedulers; the arrows show the aware reduction.

\begin{table}[t]
\caption{Per-Token Off-Chip Traffic at $B_{\mathrm{GBuf}}\!=\!512$\,MB (Qwen-1.5B)}
\label{tab:e3_qwen15b}
\centering
\begin{tabular}{lrrrr}
\toprule
Policy & PPL & Aware (MB) & Obl. (MB) & Aware $\downarrow$ \\
\midrule
INT4-uniform  & 10.452 & \textbf{124} & 392 & \textbf{3.17$\times$} \\
Mixed-4.5     & 10.279 & \textbf{213} & 454 & \textbf{2.13$\times$} \\
Mixed-5.0     & 10.091 & \textbf{289} & 516 & \textbf{1.79$\times$} \\
INT8-uniform  &  9.604 & 784 & 784 & 1.00$\times$ \\
\bottomrule
\end{tabular}
\end{table}

INT8-uniform's coincidence of aware and oblivious in
Table~\ref{tab:e3_qwen15b} is the trivial limit: the policy is already
at the worst-case bit width assumed by the oblivious scheduler, so the
conservative budget is exact.

\paragraph{Bimodal regime}
When the buffer is just large enough to hold an INT4 version of the
decoder but not an INT8 version, the aware scheduler can keep
mixed-precision models fully resident on chip and per-token off-chip
traffic drops to zero. INT8-uniform still has to stream, regardless of
scheduler. This bimodal regime is the cleanest story for an
accelerator with on-chip SRAM around the INT4-total footprint of the
model.

\subsubsection{GPU Sanity Check}
The numbers in Table~\ref{tab:e3_qwen15b} are computed analytically.
The negative mapping result of Section~\ref{sec:motivation} justifies
that approximation in the decode regime, but a hardware paper should
not rest the main result on arithmetic alone. We therefore measure the
FP16 baseline of Qwen2.5-1.5B on a single NVIDIA RTX 2000 Ada (16 GiB),
generating 256 tokens with deterministic decoding and sampling NVML
power at 20-ms intervals through a background thread.

\begin{table}[t]
\caption{GPU Sanity Profile: FP16 Qwen2.5-1.5B on a Single RTX 2000 Ada}
\label{tab:gpu_profile}
\centering
\begin{tabular}{lr}
\toprule
Metric & Value \\
\midrule
Throughput & 25.0 tokens/s \\
Wall time / 256 tokens & 10.24 s \\
Average GPU power & 68.5 W \\
Energy per token & 2.74 J \\
Power samples & 510 \\
\bottomrule
\end{tabular}
\end{table}

Table~\ref{tab:gpu_profile} is consistent with the analytic prediction:
at FP16 the per-token decoder bytes are 2.62 GB and a memory subsystem
of the relevant class caps throughput in the same neighborhood. We use
the order between this measurement and the quantized configurations as
the cross-check, not its absolute value. A controlled cross-configuration
GPU comparison is the natural follow-up.

\subsection{E3: Cross-Model Sweep}
\label{sec:e3_sweep}

\begin{figure*}[t]
\centering
\includegraphics[width=\textwidth]{../experiments/results/figs/sweep_cross_model.png}
\caption{Cross-model sweep. Solid lines: precision-aware scheduler.
Dashed: precision-oblivious INT8-budget scheduler. The aware
contribution generalizes across three model families (Qwen, Gemma,
Llama) and scales from 1.5 to 8 billion parameters. The figure shown
covers the Qwen-1.5B/Gemma-2B/Qwen-7B panels; the Llama-8B panel
mirrors the same shape and is reported in Table~\ref{tab:cross_model}.}
\label{fig:sweep_cross_model}
\end{figure*}

Figure~\ref{fig:sweep_cross_model} shows the same accuracy-versus-DRAM
relation across four models at log-scaled buffer size. Three
observations follow.

\noindent\emph{Curve shape is universal.} All four models exhibit the
same monotone-in-$B_{\mathrm{GBuf}}$ decline and the same ordering:
INT4-uniform $<$ Mixed-4.5 $<$ Mixed-5.0 $<$ INT8-uniform.

\noindent\emph{The knee sits at $B_{\mathrm{GBuf}} \approx$ INT4 total.}
For each model the knee is at the INT4-total footprint (655\,MB,
1.04\,GB, 3.25\,GB, 3.49\,GB). Below the knee, the aware-versus-oblivious
gap grows; above the knee, the aware scheduler collapses to zero.

\noindent\emph{Win factor scales with the buffer-to-decoder ratio.}
At $B_{\mathrm{GBuf}} \approx 0.4$--$0.6 \times$ INT4 total --- the
``sweet'' operating regime where the buffer holds roughly half of the
quantized decoder --- the win factors are shown in
Table~\ref{tab:cross_model}.

\begin{table}[t]
\caption{Aware-Versus-Oblivious Win at the Comparable Sweet Operating Point: the Smallest-First Knapsack and a Random-Order Knapsack Land Within One Percent of Each Other}
\label{tab:cross_model}
\centering
\small
\begin{tabular}{lrrrr}
\toprule
Model and policy   & Aware DRAM & Random DRAM & Oblivious DRAM & Oblv/Aware \\
\midrule
\multicolumn{5}{l}{\emph{Qwen-1.5B, $B_{\mathrm{GBuf}}\!=\!256$\,MB}} \\
\quad INT4-uniform  & 392.2\,MB & 386.7\,MB & 523.0\,MB & 1.33$\times$ \\
\quad Mixed-4.5     & 481.7\,MB & 480.3\,MB & 584.9\,MB & 1.21$\times$ \\
\quad Mixed-5.0     & 550.5\,MB & 550.5\,MB & 646.8\,MB & 1.18$\times$ \\
\midrule
\multicolumn{5}{l}{\emph{Gemma-2-2B, $B_{\mathrm{GBuf}}\!=\!512$\,MB}} \\
\quad INT4-uniform  & 499.0\,MB & 499.0\,MB & 775.0\,MB & \textbf{1.55$\times$} \\
\quad Mixed-4.5     & 658.2\,MB & 658.2\,MB & 902.4\,MB & 1.37$\times$ \\
\quad Mixed-5.0     & 785.6\,MB & 777.7\,MB & 998.0\,MB & 1.27$\times$ \\
\midrule
\multicolumn{5}{l}{\emph{Qwen-7B, $B_{\mathrm{GBuf}}\!=\!2$\,GB}} \\
\quad INT4-uniform  & 1.12\,GB  & 1.12\,GB  & 2.21\,GB  & \textbf{1.97$\times$} \\
\quad Mixed-4.5     & 1.60\,GB  & 1.58\,GB  & 2.31\,GB  & 1.45$\times$ \\
\quad Mixed-5.0     & 1.94\,GB  & 1.93\,GB  & 2.51\,GB  & 1.30$\times$ \\
\midrule
\multicolumn{5}{l}{\emph{Llama-8B, $B_{\mathrm{GBuf}}\!=\!2$\,GB}} \\
\quad INT4-uniform  & 1.35\,GB  & 1.34\,GB  & 2.44\,GB  & \textbf{1.80$\times$} \\
\quad Mixed-4.5     & 1.79\,GB  & 1.78\,GB  & 2.61\,GB  & 1.46$\times$ \\
\quad Mixed-5.0     & 2.23\,GB  & 2.21\,GB  & 2.97\,GB  & 1.33$\times$ \\
\bottomrule
\end{tabular}
\end{table}

Two findings follow. First, the win factor grows with the
buffer-to-decoder ratio rather than with parameter count alone: Qwen-7B
operates at $B_{\mathrm{GBuf}}/$INT4 $=$ 0.63 and reaches 1.97$\times$,
while Llama-8B operates at the same absolute
$B_{\mathrm{GBuf}}\!=\!2$\,GB but a smaller ratio of 0.59 and reaches
1.80$\times$. Per-tile byte accounting matters more as the on-chip
buffer is sized to hold a larger fraction of the four-bit decoder.

Second, the smallest-first knapsack and an uniformly-random ordering of
the same byte-aware schedule finish within roughly one percent of each
other across all sixteen (model, policy) settings; in some cases the
random order even pins marginally more bytes (e.g.,~Qwen-1.5B
INT4-uniform: 386.7\,MB versus 392.2\,MB). The implication is that
\emph{precision awareness}, not the smallest-first rule, drives almost
all of the contribution. The technique is therefore robust to solver
choice and does not depend on a particular knapsack heuristic.

\paragraph{Energy per token}
Multiplying the DRAM-byte numbers of Table~\ref{tab:cross_model} by the
Timeloop-derived constant $E_b\!=\!114.08$\,pJ/byte
(Section~\ref{sec:setup}) gives the energy per generated token at the
same sweet operating point (Table~\ref{tab:energy_per_token}).

\begin{table}[t]
\caption{Energy per Generated Token at the Sweet Operating Point (mJ/Token; Aware Versus Oblivious INT8-Budget)}
\label{tab:energy_per_token}
\centering
\small
\begin{tabular}{lrrr}
\toprule
Model and policy   & Aware & Oblivious & Aware saves \\
\midrule
\multicolumn{4}{l}{\emph{Qwen-1.5B, $B_{\mathrm{GBuf}}\!=\!256$\,MB}} \\
\quad INT4-uniform  & 44.7 & 59.7  & 14.9 \\
\quad Mixed-4.5     & 55.0 & 66.7  & 11.8 \\
\quad Mixed-5.0     & 62.8 & 73.8  & 11.0 \\
\midrule
\multicolumn{4}{l}{\emph{Gemma-2-2B, $B_{\mathrm{GBuf}}\!=\!512$\,MB}} \\
\quad INT4-uniform  & 56.9 & 88.4  & \textbf{31.5} \\
\quad Mixed-4.5     & 75.1 & 102.9 & 27.9 \\
\quad Mixed-5.0     & 89.6 & 113.9 & 24.2 \\
\midrule
\multicolumn{4}{l}{\emph{Qwen-7B, $B_{\mathrm{GBuf}}\!=\!2$\,GB}} \\
\quad INT4-uniform  & 127.8 & 251.7 & \textbf{123.9} \\
\quad Mixed-4.5     & 182.0 & 263.3 & 81.3 \\
\quad Mixed-5.0     & 220.7 & 286.6 & 65.8 \\
\midrule
\multicolumn{4}{l}{\emph{Llama-8B, $B_{\mathrm{GBuf}}\!=\!2$\,GB}} \\
\quad INT4-uniform  & 154.1 & 278.0 & \textbf{123.9} \\
\quad Mixed-4.5     & 204.3 & 298.1 & 93.8 \\
\quad Mixed-5.0     & 254.6 & 338.3 & 83.7 \\
\bottomrule
\end{tabular}
\end{table}

The absolute energy savings grow with model size: precision-aware
scheduling shaves $\sim$12--15\,mJ off every token on Qwen-1.5B but
$\sim$80--120\,mJ off every token on the 7B and 8B models, because the
DRAM-byte gap that the aware scheduler closes is larger in absolute
terms when the model is larger. At a sustained decode throughput of
100 tokens per second this corresponds to $\sim$8--12\,W of avoided
DRAM-traffic power on the 7B/8B class.

\subsection{E4: Per-Tile Hardware-Normalized Policy}
\label{sec:e4}

Section~\ref{sec:e4} reports policy shape (Table~\ref{tab:hwnorm_shape}),
GPTQ-measured perplexity (Table~\ref{tab:hwnorm_ppl}), and the residency
consequence of the resulting shape (Table~\ref{tab:hwnorm_resid}). The
accuracy result confirms the motivation for the score in
Section~\ref{sec:hwnorm}: at matched average bits per weight, the
per-tile policy beats the per-layer greedy policy on WikiText-2
perplexity by an amount that grows as the budget tightens.

The greedy per-layer score of \eqref{eq:greedy_score} ignores byte
cost. Using \eqref{eq:hwnorm} at per-projection granularity (196 tiles
for Qwen-1.5B), Table~\ref{tab:hwnorm_shape} contrasts the resulting
policy shape against the per-layer greedy at the same average bits per
weight.

\begin{table}[t]
\caption{Policy Shape: Per-Layer Greedy Versus Per-Tile Hardware-Normalized at $\bar B = 4.5$}
\label{tab:hwnorm_shape}
\centering
\begin{tabular}{lrr}
\toprule
Projection & Greedy (out of 28) & Hwnorm (out of 28) \\
\midrule
\texttt{q\_proj}    & 4 at 8-bit & 7 at 8-bit \\
\texttt{k\_proj}    & 4          & \textbf{17}        \\
\texttt{v\_proj}    & 4          & \textbf{25 (89\%)} \\
\texttt{o\_proj}    & 4          & 15         \\
\texttt{gate\_proj} & 4          & \textbf{0 (all 4-bit)} \\
\texttt{up\_proj}   & 4          & 2          \\
\texttt{down\_proj} & 4          & 5          \\
\midrule
Total 8-bit tiles & 28 & \textbf{71} \\
Achieved BPW      & 4.571 & \textbf{4.503} \\
\bottomrule
\end{tabular}
\end{table}

The hardware-normalized policy concentrates the eight-bit budget in
attention's KV projections, which are the cheapest to promote
(approximately four to six times smaller than MLP projections) and
remain sensitive. MLP \texttt{gate\_proj} stays entirely at INT4
because it is 6.89~MB per layer at INT4 and any single promotion to
eight-bit consumes 6.89~MB of budget --- too expensive against its
$\Delta\mathrm{PPL}$ contribution.

\begin{table}[t]
\caption{Real GPTQ Perplexity at Matched Average Bit-Width: Per-Tile Hardware-Normalized Beats Per-Layer Greedy in Seven of Eight Settings Across Four Models}
\label{tab:hwnorm_ppl}
\centering
\small
\begin{tabular}{llrrr}
\toprule
Model & Budget & Greedy ppl & Hwnorm ppl & $\Delta$ ppl \\
\midrule
\multirow{4}{*}{Qwen-1.5B}  & $\bar B = 4.5$ & 10.279 & \textbf{10.117} & $-0.162$ \\
                            & $\bar B = 5.0$ & 10.091 & \textbf{10.025} & $-0.066$ \\
                            & $\bar B = 5.5$ & ---    & 9.912           & ---      \\
                            & $\bar B = 6.0$ & ---    & 9.838           & ---      \\
\midrule
\multirow{2}{*}{Gemma-2-2B} & $\bar B = 4.5$ & 14.723 & \textbf{14.663} & $-0.060$ \\
                            & $\bar B = 5.0$ & 14.387 & 14.406          & $+0.019$ \\
\midrule
\multirow{2}{*}{Qwen-7B}    & $\bar B = 4.5$ & 7.724  & \textbf{7.676}  & $-0.048$ \\
                            & $\bar B = 5.0$ & 7.661  & \textbf{7.581}  & $-0.080$ \\
\midrule
\multirow{2}{*}{Llama-8B}   & $\bar B = 4.5$ & 7.561  & \textbf{7.532}  & $-0.029$ \\
                            & $\bar B = 5.0$ & 7.514  & \textbf{7.467}  & $-0.047$ \\
\bottomrule
\multicolumn{5}{l}{\footnotesize Greedy ppl for Qwen-1.5B at $\bar B \in \{5.5, 6.0\}$ was not run in E1.}
\end{tabular}
\end{table}

Across the four models, the per-tile hardware-normalized policy
reduces WikiText-2 perplexity at matched average bit-width in seven of
the eight settings reported in Table~\ref{tab:hwnorm_ppl}. The lone
exception is Gemma-2-2B at $\bar B\!=\!5.0$, where the gap is $+$0.019
--- within the magnitude we attribute to GPTQ calibration noise at 32
samples. The win is concentrated at the tightest budget ($\bar B\!=\!4.5$),
which is consistent with the intuition that protecting key-value
projections matters most when the eight-bit promotion budget is small.

\begin{table}[t]
\caption{Residency Consequence at $B_{\mathrm{GBuf}}\!=\!512$\,MB Is Modest}
\label{tab:hwnorm_resid}
\centering
\begin{tabular}{lrrrr}
\toprule
Budget & Greedy BPW & Hwnorm BPW & Greedy DRAM & Hwnorm DRAM \\
\midrule
$\bar B = 4.5$ & 4.571 & 4.503 & 213.3\,MB & \textbf{206.4\,MB} \\
$\bar B = 5.0$ & 5.000 & 5.025 & 289.0\,MB & 289.0\,MB \\
\bottomrule
\end{tabular}
\end{table}

The residency win is small because the knapsack picks tiles by absolute
size; it is unaware of which logical role each tile plays
(Table~\ref{tab:hwnorm_resid}). The accuracy side is where we expect
the real win: protecting 89\% of \texttt{v\_proj} at eight-bit while
keeping all \texttt{gate\_proj} at four-bit yields a lower perplexity
at iso-budget than the per-layer policy, as confirmed by the GPTQ
results of Table~\ref{tab:hwnorm_ppl}.

\section{Discussion}
\label{sec:discussion}

\subsection{What the Negative Mapping Result Really Means}
Section~\ref{sec:motivation} reports that for a 16$\times$16
Eyeriss-class architecture at $N\!=\!1$ decode, the mapper finds
equivalent-cost mappings for INT4 and INT8 versions of the same matrix
multiply. This is not a claim that mapping search is useless for LLM
accelerators; only that, in the regime where total weight bytes are
much larger than on-chip capacity and each weight is touched once per
token, the mapping decision space collapses onto a single
byte-bound point. In two complementary regimes we expect mapping to
recover relevance: prefill ($N \gg 1$), where weights are reused across
tokens within the forward pass, and architectures with comparable
on-chip and off-chip ratios, where buffer-pressure effects become the
dominant variable.

\subsection{Limitations and Follow-Ups}
\label{sec:limitations}
First, our energy estimate of Table~\ref{tab:energy_per_token} is a
DRAM-traffic lower bound derived from the per-byte energy harvested by
Timeloop; it omits PE-array idle power and on-chip-buffer refresh,
which are common across schedulers but inflate the absolute mJ/token
number. A full Accelergy run that integrates idle and leakage power
would tighten the absolute energy, though not the aware-versus-oblivious
ratios reported here. Second, the zero-one knapsack
model is conservative: it pins or streams each tile entirely. A real
scheduler could prefetch the next layer's weights while the current
layer computes; modeling that would reduce the oblivious baseline's
overhead. Third, pinning is valuable across many tokens; for a
single-token generation the residency win is zero. Our results
implicitly assume sustained decoding. Fourth, the per-tile
hardware-normalized policy is validated on Qwen-1.5B with four budgets
and on Gemma-2-2B, Qwen-7B, and Llama-8B with two budgets each
(Table~\ref{tab:hwnorm_ppl}); replicating it on additional model
families and at the loosest budgets would round out the picture.
Fifth, downstream-task evaluation is reported on a 1000-sample
HellaSwag subset (Table~\ref{tab:hellaswag}); a larger subset and
additional tasks such as ARC-Challenge and MMLU would sharpen the
sub-2-percent gaps that fall inside our standard-error band at the
present sample size.

\subsection{Composability With Mapping-Layer Work}
Power-aware mapping such as MAGNETO (anonymized) and LLM-specific
accelerator design such as GATHER (anonymized) sit at the mapping
layer of Fig.~\ref{fig:stack}. Our work sits one layer higher and
composes with both: given a fixed policy and a fixed mapping, the
scheduler decides which tiles benefit from on-chip residency.

\section{Related Work}
\label{sec:related}

\subsection{Layer-Wise Mixed-Precision Policy Search}
HAWQ \cite{dong2019hawq} and its eigenvalue variant
\cite{dong2020hawqv2} introduced second-order sensitivity scoring for
per-layer bit-width assignment. BRECQ \cite{li2021brecq} formulated
block-wise reconstruction error as a more practical signal, and
OmniQuant \cite{shao2024omniquant} re-parameterized the search to make
the budget a learnable variable. LLM.int8() \cite{dettmers2022llmint8}
argued for outlier-aware INT8 inference. All four sit at the algorithm
layer of Fig.~\ref{fig:stack}; none addresses how the resulting policy
interacts with the deployment-time residency budget.

\subsection{LLM Quantization at Deployment}
GPTQ \cite{frantar2023gptq} performs Hessian-aware column-wise
rounding; we use it as our quantizer of record. AWQ
\cite{lin2024awq} preserves activation-salient channels;
SmoothQuant \cite{xiao2023smoothquant} shifts the
activation-versus-weight quantization difficulty. Marlin
\cite{frantar2024marlin} provides a fast four-bit decode kernel that
turns the byte saving of mixed precision into latency savings on
modern GPUs. These produce kernels; scheduling decisions about which
weights stay on chip are outside their scope.

\subsection{LLM Accelerator Dataflow}
Sanger \cite{lu2021sanger} and Ant \cite{guo2022ant} are representative
LLM accelerator designs with custom dataflow. Our laboratory's prior
power-aware mapping and LLM-specific accelerator (both anonymized for
review) sit in the same space. Each contributes at the mapping layer
of Fig.~\ref{fig:stack}; they assume a fixed precision plan as input
and search the tile and dataflow space. Our work composes with all of
them: a Sanger-class mapper produces the per-shape mapping that
determines cycles per token per tile; our scheduler decides which of
those tiles avoid DRAM round trips.

\subsection{Memory Residency and Cache Scheduling for LLM Serving}
SARATHI \cite{agrawal2024sarathi} schedules prefill and decode batches
to improve GPU utilization but operates at the request level.
DistServe \cite{zhong2024distserve} decouples prefill and decode
across machines. FlexGen \cite{sheng2023flexgen} tiers weight storage
across GPU, CPU, and SSD for offload-friendly inference.
ChunkAttention \cite{ye2024chunkattention} re-orders KV-cache access
for sequence reuse. These works address the residency of the
key-value cache or of runtime state. Our work addresses the residency
of weights under a mixed-precision policy, a different optimization
variable.

\subsection{Knapsack-Based Scheduling for Inference}
Tetris \cite{gao2017tetris} uses integer-linear-programming over
operators for three-dimensional DNN partitioning; TVM/Ansor
\cite{zheng2020ansor} is search-based but at the kernel-fusion
granularity. None of these are precision-aware: the operand bit width
is fixed upfront. Our knapsack instance is defined by the per-tile
bytes, which are constructed by the upstream policy; that is the
precise coupling that the scheduling layer must surface.

\section{Conclusion}
\label{sec:conclusion}
We isolate a system-level question that the mixed-precision quantization
literature does not address: given a fixed policy, which weight tiles
do we keep on-chip? We formulate the problem as a greedy zero-one
knapsack over per-tile bytes, show that a precision-aware scheduler
reduces per-token off-chip traffic by 1.33$\times$--1.97$\times$ across
three LLM families (Qwen, Gemma, Llama) and four checkpoints from 1.5 to
8 billion parameters at a comparable operating point, and up to
3.17$\times$ in the knee regime for the 1.5B-parameter model. The
contribution generalizes with a factor that grows with the
buffer-to-decoder ratio. We release the full reproducibility trail, including an
analytic-to-measurement cross-check via a GPU profile of the
unquantized baseline. The result puts a name and a quantification on a
layer of the stack between the quantization algorithm and the dataflow
mapper, and points to prefetch-aware scheduling as the natural
follow-up. As a method extension, the per-tile
hardware-normalized policy reduces WikiText-2 perplexity in seven of
eight settings across four models at matched average bit-width
(Section~\ref{sec:e4}, Table~\ref{tab:hwnorm_ppl}), with the largest
gain ($-0.162$ at $\bar B = 4.5$ on Qwen-1.5B) concentrated at the
tightest budget where protecting key-value projections matters most.
We further demonstrate that the win is driven by precision awareness
itself: replacing the smallest-first knapsack with a uniformly-random
order on the same byte-aware budget shifts DRAM traffic by less than
one percent, so the result is robust to solver choice.

\section*{Acknowledgment}
The authors thank the members of our laboratory for early discussion
on the framing.

\noindent\textbf{Use of generative AI.} Portions of the text in
Sections~\ref{sec:introduction}, \ref{sec:method}, \ref{sec:results},
\ref{sec:related}, and \ref{sec:conclusion} were drafted with the
assistance of Anthropic Claude \cite{anthropic2024claude}. All
technical claims, equations, and the experimental design are the
authors' own; the AI's role was limited to language polishing,
structural editing, and the suggestion of alternative phrasings, after
which the authors verified every quantitative statement against the
experimental artifacts. No part of the experimental design,
implementation, or numerical results was generated by the AI without
authorial verification.

\begin{thebibliography}{99}

\bibitem{dong2019hawq}
Z.~Dong, Z.~Yao, A.~Gholami, M.~W. Mahoney, and K.~Keutzer,
``HAWQ: Hessian AWare Quantization of neural networks with mixed-precision,''
in \emph{Proc. IEEE/CVF Int. Conf. Comput. Vis. (ICCV)}, Seoul, South Korea, Oct. 2019, pp. 293--302.

\bibitem{dong2020hawqv2}
Z.~Dong, Z.~Yao, D.~Arfeen, A.~Gholami, M.~W. Mahoney, and K.~Keutzer,
``HAWQ-V2: Hessian aware trace-weighted quantization of neural networks,''
in \emph{Adv. Neural Inf. Process. Syst. (NeurIPS)}, vol. 33, Dec. 2020, pp. 18518--18529.

\bibitem{li2021brecq}
Y.~Li, R.~Gong, X.~Tan, Y.~Yang, P.~Hu, Q.~Zhang, F.~Yu, W.~Wang, and S.~Gu,
``BRECQ: Pushing the limit of post-training quantization by block reconstruction,''
in \emph{Proc. Int. Conf. Learn. Represent. (ICLR)}, May 2021.

\bibitem{shao2024omniquant}
W.~Shao, M.~Chen, Z.~Zhang, P.~Xu, L.~Zhao, Z.~Li, K.~Zhang, P.~Gao, Y.~Qiao, and P.~Luo,
``OmniQuant: Omnidirectionally calibrated quantization for large language models,''
in \emph{Proc. Int. Conf. Learn. Represent. (ICLR)}, May 2024.

\bibitem{dettmers2022llmint8}
T.~Dettmers, M.~Lewis, Y.~Belkada, and L.~Zettlemoyer,
``LLM.int8(): 8-bit matrix multiplication for transformers at scale,''
in \emph{Adv. Neural Inf. Process. Syst. (NeurIPS)}, vol. 35, Dec. 2022, pp. 30318--30332.

\bibitem{frantar2023gptq}
E.~Frantar, S.~Ashkboos, T.~Hoefler, and D.~Alistarh,
``GPTQ: Accurate post-training quantization for generative pre-trained transformers,''
in \emph{Proc. Int. Conf. Learn. Represent. (ICLR)}, May 2023.

\bibitem{lin2024awq}
J.~Lin, J.~Tang, H.~Tang, S.~Yang, X.~Dang, C.~Gan, and S.~Han,
``AWQ: Activation-aware weight quantization for LLM compression and acceleration,''
in \emph{Proc. Mach. Learn. Syst. (MLSys)}, May 2024.

\bibitem{xiao2023smoothquant}
G.~Xiao, J.~Lin, M.~Seznec, H.~Wu, J.~Demouth, and S.~Han,
``SmoothQuant: Accurate and efficient post-training quantization for large language models,''
in \emph{Proc. Int. Conf. Mach. Learn. (ICML)}, vol. 202, Jul. 2023, pp. 38087--38099.

\bibitem{frantar2024marlin}
E.~Frantar, R.~Castro, J.~Chen, T.~Hoefler, and D.~Alistarh,
``Marlin: Mixed-precision auto-regressive parallel inference on large language models,''
arXiv preprint arXiv:2408.11743, 2024.

\bibitem{parashar2019timeloop}
A.~Parashar, P.~Raina, Y.~S. Shao, Y.-H. Chen, V.~A. Ying, A.~Mukkara, R.~Venkatesan, B.~Khailany, S.~W. Keckler, and J.~Emer,
``Timeloop: A systematic approach to DNN accelerator evaluation,''
in \emph{Proc. IEEE Int. Symp. Perform. Anal. Syst. Softw. (ISPASS)}, Madison, WI, USA, Mar. 2019, pp. 304--315.

\bibitem{chen2017eyeriss}
Y.-H. Chen, T.~Krishna, J.~S. Emer, and V.~Sze,
``Eyeriss: An energy-efficient reconfigurable accelerator for deep convolutional neural networks,''
\emph{IEEE J. Solid-State Circuits}, vol. 52, no. 1, pp. 127--138, Jan. 2017.

\bibitem{lu2021sanger}
L.~Lu, Y.~Jin, H.~Bi, Z.~Luo, P.~Li, T.~Wang, and Y.~Liang,
``Sanger: A co-design framework for enabling sparse attention using reconfigurable architecture,''
in \emph{Proc. 54th Annu. IEEE/ACM Int. Symp. Microarchitecture (MICRO)}, Athens, Greece, Oct. 2021, pp. 977--991.

\bibitem{guo2022ant}
C.~Guo, C.~Zhang, J.~Leng, Z.~Liu, F.~Yang, Y.~Liu, M.~Guo, and Y.~Zhu,
``ANT: Exploiting adaptive numerical data type for low-bit deep neural network quantization,''
in \emph{Proc. 55th Annu. IEEE/ACM Int. Symp. Microarchitecture (MICRO)}, Chicago, IL, USA, Oct. 2022, pp. 1414--1428.

\bibitem{sheng2023flexgen}
Y.~Sheng, L.~Zheng, B.~Yuan, Z.~Li, M.~Ryabinin, B.~Chen, P.~Liang, C.~R\'e, I.~Stoica, and C.~Zhang,
``FlexGen: High-throughput generative inference of large language models with a single GPU,''
in \emph{Proc. Int. Conf. Mach. Learn. (ICML)}, vol. 202, Jul. 2023, pp. 31094--31116.

\bibitem{agrawal2024sarathi}
A.~Agrawal, A.~Panwar, J.~Mohan, N.~Kwatra, B.~S. Gulavani, and R.~Ramjee,
``SARATHI: Efficient LLM inference by piggybacking decodes with chunked prefills,''
in \emph{Proc. Mach. Learn. Syst. (MLSys)}, May 2024.

\bibitem{zhong2024distserve}
Y.~Zhong, S.~Liu, J.~Chen, J.~Hu, Y.~Zhu, X.~Liu, X.~Jin, and H.~Zhang,
``DistServe: Disaggregating prefill and decoding for goodput-optimized large language model serving,''
in \emph{Proc. USENIX Symp. Oper. Syst. Des. Implement. (OSDI)}, Santa Clara, CA, USA, Jul. 2024, pp. 193--210.

\bibitem{ye2024chunkattention}
L.~Ye, Z.~Tao, Y.~Huang, and Y.~Li,
``ChunkAttention: Efficient self-attention with prefix-aware KV cache and two-phase partition,''
in \emph{Proc. Annu. Meet. Assoc. Comput. Linguistics (ACL)}, Aug. 2024.

\bibitem{gao2017tetris}
M.~Gao, J.~Pu, X.~Yang, M.~Horowitz, and C.~Kozyrakis,
``TETRIS: Scalable and efficient neural network acceleration with 3D memory,''
in \emph{Proc. 22nd Int. Conf. Architectural Support Program. Lang. Oper. Syst. (ASPLOS)}, Xi'an, China, Apr. 2017, pp. 751--764.

\bibitem{zheng2020ansor}
L.~Zheng, C.~Jia, M.~Sun, Z.~Wu, C.~H. Yu, A.~Haj-Ali, Y.~Wang, J.~Yang, D.~Zhuo, K.~Sen, J.~E. Gonzalez, and I.~Stoica,
``Ansor: Generating high-performance tensor programs for deep learning,''
in \emph{Proc. USENIX Symp. Oper. Syst. Des. Implement. (OSDI)}, Nov. 2020, pp. 863--879.

\bibitem{accelerate}
S.~Gugger, L.~Debut, T.~Wolf, P.~Schmid, Z.~Mueller, S.~Mangrulkar, M.~Sun, and B.~Bossan,
``Accelerate: Training and inference at scale made simple, efficient and adaptable,''
\emph{Hugging Face}, GitHub: huggingface/accelerate, 2022.

\bibitem{anthropic2024claude}
Anthropic,
``Claude,'' Anthropic, San Francisco, CA, USA, 2024. [Online].
Available: \url{https://www.anthropic.com/claude}

\end{thebibliography}

\EOD

\begin{IEEEbiography}{Author One} (placeholder) is with the Department
of Computer Science and Engineering, Ewha Womans University, Seoul,
Republic of Korea. The author's research interests include machine
learning systems, accelerator architecture, and post-training
quantization for large language models. A biography paragraph, ORCID,
and an author photograph will be provided in the camera-ready version.
\end{IEEEbiography}

\begin{IEEEbiography}{Author Two} (placeholder) is with the Department
of Computer Science and Engineering, Ewha Womans University, Seoul,
Republic of Korea. The author's interests include accelerator mapping,
power-aware dataflow, and hardware--software co-design for deep neural
networks. A biography paragraph, ORCID, and an author photograph will
be provided in the camera-ready version.
\end{IEEEbiography}

\begin{IEEEbiography}{Author Three} (placeholder) is with the Department
of Computer Science and Engineering, Ewha Womans University, Seoul,
Republic of Korea. The author's interests include LLM-specific
accelerator design, on-chip memory hierarchies, and inference-time
scheduling. A biography paragraph, ORCID, and an author photograph will
be provided in the camera-ready version.
\end{IEEEbiography}

\end{document}
